% Generated draft: requires mathematical and source review.
\documentclass[11pt]{article}
\usepackage[margin=25mm]{geometry}
\usepackage{fontspec}
\setmainfont{Latin Modern Roman}
\usepackage{amsmath,amssymb,mathtools,unicode-math}
\setmathfont{Latin Modern Math}
\usepackage{enumitem,needspace,xurl,hyperref}
\setlength{\parindent}{0pt}
\setlength{\parskip}{5pt}
\setlength{\emergencystretch}{3em}
\newcommand{\N}{\mathbb{N}}
\newcommand{\Z}{\mathbb{Z}}
\newcommand{\Q}{\mathbb{Q}}
\newcommand{\R}{\mathbb{R}}
\newcommand{\C}{\mathbb{C}}
\newcommand{\F}{\mathbb{F}}
\newcommand{\A}{\mathbb{A}}
\newcommand{\PP}{\mathbb{P}}
\newcommand{\E}{\mathbb{E}}
\newcommand{\eps}{\varepsilon}
\newcommand{\dd}{\,\mathrm{d}}
\begin{document}
\section*{UnsolvedMath ID: 30006568}
\textbf{Source problem code:} catalogue entry
\subsection{Definitions and mathematical statement}
Let \(A\) be a matrix satisfying the hypotheses of the limiting quantum isometry construction referred to in the source. Write \(\Sigma_A\) for its path space and \(G_A^\infty\) for the limiting quantum isometry group. The supplied excerpt does not specify the admissible matrices, the convention for paths, or the construction of \(G_A^\infty\); these objects must therefore be understood conditionally on the missing source definitions. In particular, the assertion that \(\Sigma_A\) is a Cantor set requires the source's hypotheses and is not asserted here for arbitrary matrices.

Denote the induced continuous compact quantum group action by its coaction
\[
\alpha_A:C(\Sigma_A)\longrightarrow C(\Sigma_A)\otimes_{\min}C(G_A^\infty),
\]
where \(C(\Sigma_A)\) is the algebra of continuous complex-valued functions on the compact path space and \(C(G_A^\infty)\) is the function algebra of the compact quantum group. The action is called \emph{ergodic} when its fixed-point algebra consists exactly of the constant functions:
\[
C(\Sigma_A)^{\alpha_A}
:=\{f\in C(\Sigma_A):\alpha_A(f)=f\otimes 1\}
=\mathbb C1.
\]
Let \(\mathbf 1\) denote the full-shift matrix of the relevant size, all of whose entries are one; the source's notation \(A=1\) has this meaning.

\textbf{Problem.} Establish ergodicity of the source's limiting quantum isometry action for some admissible matrices \(A\ne\mathbf 1\). Thus, identify matrices beyond the full-shift case for which the source construction is defined and prove
\[
\forall f\in C(\Sigma_A),\qquad
\alpha_A(f)=f\otimes1
\ \Longrightarrow\ 
f\in\mathbb C1.
\]
The original question asks for some cases beyond the full shift; it does not specify a universal assertion for every admissible matrix.

\textbf{Supplied context and unresolved qualifications.} The excerpt attributes to its Theorem 3 the preservation of \(\Sigma_A\) by \(G_A^\infty\). For the full-shift case it specifically states that a certain quotient \(H_1\) of \(G_{\mathbf 1}^\infty\), itself a nonclassical compact quantum group, acts faithfully and ergodically on \(\Sigma_{\mathbf 1}\). The quotient, its action, and its relationship to the action denoted by \(\alpha_A\) above are not defined in the supplied record. Faithfulness is part of that reported full-shift result; the requested extension asks for ergodicity.
\subsection{Short English statement}
Find cases beyond the all-ones full-shift matrix in which the limiting quantum isometry action on the associated path space has only constant invariant continuous functions. The admissible matrices and the limiting action must be recovered from the source before this conditional formulation becomes a fully specified problem. The supplied record labels the problem open; that label is not an independently verified conclusion.
\subsection{Sources}
\begin{itemize}
\item Supplied local UnsolvedMath registry record, permanent ID \(30006568\), title \emph{Ergodic Limiting Quantum Isometry Actions}, exact source problem code \texttt{catalogue entry}. No individual proposer is identified. The registry records proposed year 2026, status \texttt{open}, and statement status \texttt{corrected\_verified}; its statement-verification field says the reconstruction was checked against the official report. These are supplied registry assertions, not independent verification here.
\item \emph{Noncommutative Harmonic Analysis and Quantum Information} (2026), Oberwolfach Reports, DOI \texttt{10.4171/owr/2026/4}, pp.\ 217--275. Supplied source and literature-source link, labelled ``Oberwolfach Report 4/2026'': \url{https://doi.org/10.4171/owr/2026/4}. The supplied background describes these reports as documenting mathematical workshops and their open-problem discussions. An authoritative report citation is supplied, but the excerpt omits the preceding construction and its hypotheses.
\item Literature assessment dated \textbf{2026-08-22}: the supplied curation labels the problem open, states that the report leaves ergodicity beyond the full-shift case unanswered, and reports that no later result establishing the general claim was located. This is supplied dated curation, not an independently verified literature assessment; its time-sensitive status and cited claims require independent checking.
\item The excerpt cites \textnormal{[GJ18]} for a comparison with quantum isometry groups of classical smooth manifolds. No authors, title, bibliographic details, or link for that reference are supplied, so its identity and the precise hypotheses of the comparison remain unresolved.
\end{itemize}
\end{document}
